\documentclass[aspectratio=169,11pt]{beamer}

\usepackage[utf8]{inputenc}
\usepackage[T1]{fontenc}
\usepackage[brazil]{babel}
\usepackage{amsmath}
\usepackage{booktabs}
\usepackage{array}
\usepackage{xcolor}

\usetheme{Madrid}
\usecolortheme{default}

\definecolor{uftblue}{RGB}{0,76,153}
\definecolor{uftgreen}{RGB}{0,128,96}
\definecolor{uftred}{RGB}{180,55,55}
\definecolor{softgray}{RGB}{245,247,250}

\setbeamercolor{structure}{fg=uftblue}
\setbeamercolor{frametitle}{fg=white,bg=uftblue}
\setbeamercolor{title}{fg=white,bg=uftblue}
\setbeamercolor{block title}{fg=white,bg=uftblue}
\setbeamercolor{block body}{bg=softgray}
\setbeamertemplate{navigation symbols}{}
\setbeamertemplate{itemize item}{\color{uftblue}$\blacktriangleright$}
\setbeamertemplate{itemize subitem}{\color{uftgreen}$\bullet$}

\title[Modelo ZIP no ANAREDE]{Análise de Fluxo de Potência com Modelagem ZIP de Carga}
\subtitle{Interpretação dos resultados no ANAREDE com base na teoria de SEP}
\author{Sistema IEEE 39 barras}
\institute{Sistemas Elétricos de Potência}
\date{\today}

\newcommand{\casobase}{Caso Base}
\newcommand{\casozip}{Caso ZIP}
\newcommand{\casozipcento}{Caso ZIP 110\%}
\newcommand{\pu}{p.u.}

\begin{document}

\begin{frame}
  \titlepage
\end{frame}

\begin{frame}{Roteiro}
  \tableofcontents
\end{frame}

\section{Objetivo e Contexto}

\begin{frame}{Objetivo da Apresentação}
  \begin{block}{Objetivo geral}
    Apresentar como a modelagem ZIP de carga altera os resultados de fluxo de potência no ANAREDE.
  \end{block}

  \vspace{0.3cm}
  A análise compara três condições de operação:

  \begin{itemize}
    \item \textbf{\casobase}: caso de referência, com cargas representadas no modelo convencional.
    \item \textbf{\casozip}: caso com cargas representadas por modelo dependente da tensão.
    \item \textbf{\casozipcento}: cenário adicional com alteração no carregamento ZIP.
  \end{itemize}

  \vspace{0.2cm}
  A comparação será feita por tensões, ângulos, fluxos ativos e reativos, geração, carga e perdas.
\end{frame}

\begin{frame}{Base Teórica Utilizada}
  A fundamentação vem do capítulo de fluxo de potência e da seção de modelagem de carga do material da disciplina:

  \begin{block}{Material de referência}
    \texttt{docs/material\_agrupado\_sep2 (1).pdf}
  \end{block}

  O material apresenta que o fluxo de potência fornece:

  \begin{itemize}
    \item tensões e ângulos nas barras;
    \item potências ativa e reativa nos ramos;
    \item perdas elétricas;
    \item condição de operação do sistema em regime permanente.
  \end{itemize}
\end{frame}

\begin{frame}{Por que Fazer Fluxo de Potência?}
  O fluxo de potência é o primeiro estudo operacional de um SEP porque estabelece o balanço entre geração, carga e perdas.

  \vspace{0.25cm}
  Com ele é possível responder:

  \begin{itemize}
    \item as tensões estão dentro de uma faixa aceitável?
    \item quais linhas e transformadores estão mais carregados?
    \item qual é o sentido do fluxo de potência?
    \item quais são as perdas elétricas do sistema?
    \item como o sistema reage a mudanças na carga?
  \end{itemize}

  \vspace{0.25cm}
  Nesta apresentação, o foco é a resposta do sistema quando a carga deixa de ser tratada como puramente constante.
\end{frame}

\section{Teoria do Fluxo de Potência}

\begin{frame}{Grandezas Calculadas no Fluxo}
  Em uma rede elétrica, o fluxo de potência calcula as variáveis de regime permanente.

  \begin{columns}[T]
    \begin{column}{0.48\textwidth}
      \begin{block}{Barras}
        \begin{itemize}
          \item módulo da tensão: $|V|$;
          \item ângulo da tensão: $\theta$;
          \item potência gerada: $P_G$, $Q_G$;
          \item potência consumida: $P_L$, $Q_L$.
        \end{itemize}
      \end{block}
    \end{column}

    \begin{column}{0.48\textwidth}
      \begin{block}{Ramos}
        \begin{itemize}
          \item fluxo ativo: $P_{km}$;
          \item fluxo reativo: $Q_{km}$;
          \item perdas ativas;
          \item redistribuição de carregamento.
        \end{itemize}
      \end{block}
    \end{column}
  \end{columns}
\end{frame}

\begin{frame}{Natureza Não Linear do Problema}
  O fluxo de potência é um problema não linear porque potência depende do produto entre tensão e corrente:

  \[
    S = P + jQ = V I^{*}
  \]

  A corrente depende da rede:

  \[
    I = Y_{\text{bus}} V
  \]

  Logo, as potências injetadas nas barras dependem das magnitudes e dos ângulos das tensões.

  \begin{block}{Consequência prática}
    Mudanças na carga, nos taps, na compensação reativa ou na topologia alteram as tensões, os ângulos, os fluxos e as perdas.
  \end{block}
\end{frame}

\begin{frame}{Tipos de Barra no Fluxo}
  No estudo de fluxo de potência, as barras são classificadas conforme as grandezas conhecidas.

  \begin{table}
    \centering
    \begin{tabular}{lll}
      \toprule
      Tipo de barra & Dados conhecidos & Dados calculados \\
      \midrule
      PQ & $P$ e $Q$ & $|V|$ e $\theta$ \\
      PV & $P$ e $|V|$ & $Q$ e $\theta$ \\
      Referência & $|V|$ e $\theta$ & $P$ e $Q$ \\
      \bottomrule
    \end{tabular}
  \end{table}

  \vspace{0.2cm}
  No ANAREDE, essa classificação permite resolver a rede e determinar o estado operativo do sistema.
\end{frame}

\section{Modelagem ZIP de Carga}

\begin{frame}{Motivação para o Modelo ZIP}
  A carga real de uma barra não depende apenas do valor nominal em MW e MVAr.

  \vspace{0.25cm}
  Segundo a teoria da modelagem de carga, a demanda também depende da tensão aplicada:

  \[
    P = f(P_0,V)
    \qquad
    Q = f(Q_0,V)
  \]

  \begin{block}{Ideia central}
    Quando a tensão se afasta de $1{,}0$ \pu, a potência consumida por certos tipos de carga também muda.
  \end{block}

  Isso é relevante porque quedas ou elevações de tensão podem alterar o carregamento efetivo do sistema.
\end{frame}

\begin{frame}{As Três Parcelas do Modelo ZIP}
  O nome ZIP vem das três formas de comportamento da carga:

  \begin{table}
    \centering
    \begin{tabular}{lll}
      \toprule
      Parcela & Comportamento & Relação com a tensão \\
      \midrule
      Z & Impedância constante & $P,Q \propto V^2$ \\
      I & Corrente constante & $P,Q \propto V$ \\
      P & Potência constante & $P,Q =$ constante \\
      \bottomrule
    \end{tabular}
  \end{table}

  \vspace{0.2cm}
  Assim, uma mesma barra pode ter parte da carga sensível ao quadrado da tensão, parte sensível linearmente e parte independente da tensão.
\end{frame}

\begin{frame}{Equação da Potência Ativa no ZIP}
  Para a potência ativa, o modelo ZIP pode ser escrito como:

  \[
    P = P_0 \left[
      k_{pp}
      + k_{pi}\left(\frac{V}{V_0}\right)
      + k_{pz}\left(\frac{V}{V_0}\right)^2
    \right]
  \]

  com:

  \[
    k_{pz} + k_{pi} + k_{pp} = 1
  \]

  \begin{block}{Interpretação}
    Se a tensão cai, as parcelas Z e I tendem a reduzir a potência ativa efetiva. A parcela P mantém a potência constante.
  \end{block}
\end{frame}

\begin{frame}{Equação da Potência Reativa no ZIP}
  A mesma lógica é aplicada à potência reativa:

  \[
    Q = Q_0 \left[
      k_{qp}
      + k_{qi}\left(\frac{V}{V_0}\right)
      + k_{qz}\left(\frac{V}{V_0}\right)^2
    \right]
  \]

  com:

  \[
    k_{qz} + k_{qi} + k_{qp} = 1
  \]

  \begin{block}{Importância no SEP}
    A potência reativa influencia diretamente o perfil de tensão. Por isso, alterações em $Q$ podem aparecer como quedas ou elevações de tensão nas barras.
  \end{block}
\end{frame}

\begin{frame}{Efeito Esperado da Modelagem ZIP}
  Com o modelo ZIP, espera-se que:

  \begin{itemize}
    \item barras com tensão menor que $1{,}0$ \pu{} tenham redução nas parcelas Z e I da carga;
    \item barras com tensão maior que $1{,}0$ \pu{} tenham aumento nas parcelas Z e I da carga;
    \item a distribuição dos fluxos mude, pois as cargas efetivas mudam por barra;
    \item as perdas mudem, porque dependem da circulação de corrente nos ramos.
  \end{itemize}

  \begin{block}{Ponto-chave}
    A modelagem ZIP não é apenas uma alteração de entrada: ela muda o comportamento elétrico do sistema resolvido pelo fluxo de potência.
  \end{block}
\end{frame}

\section{Metodologia}

\begin{frame}{Sistema e Ferramenta}
  A simulação foi realizada no ANAREDE, considerando o sistema IEEE 39 barras.

  \begin{columns}[T]
    \begin{column}{0.48\textwidth}
      \begin{block}{Sistema}
        \begin{itemize}
          \item 39 barras;
          \item ramos de transmissão;
          \item barras de carga;
          \item barras de geração;
          \item barra de referência.
        \end{itemize}
      \end{block}
    \end{column}

    \begin{column}{0.48\textwidth}
      \begin{block}{Cenários}
        \begin{itemize}
          \item \casobase;
          \item \casozip;
          \item \casozipcento.
        \end{itemize}
      \end{block}
    \end{column}
  \end{columns}
\end{frame}

\begin{frame}{Sequência da Análise}
  A metodologia usada na comparação foi:

  \begin{enumerate}
    \item Executar o fluxo de potência do \casobase.
    \item Executar o fluxo de potência do \casozip.
    \item Executar o fluxo de potência do \casozipcento.
    \item Extrair tensões, ângulos, cargas, gerações, fluxos e perdas.
    \item Comparar:
      \begin{itemize}
        \item \casobase{} $\times$ \casozip;
        \item \casozip{} $\times$ \casozipcento.
      \end{itemize}
  \end{enumerate}

  \vspace{0.15cm}
  A análise considera variações em \pu, graus, MW e MVAr.
\end{frame}

\begin{frame}{Observação Sobre o Aviso do Diagrama}
  Ao carregar alguns diagramas dos casos ZIP, apareceu aviso de diferença no grupo base:

  \begin{block}{Aviso do ANAREDE}
    O grupo base do diagrama indicava 345 kV, enquanto o caso indicava 0 kV para o mesmo identificador.
  \end{block}

  Esse aviso está associado à importação gráfica do diagrama.

  \begin{itemize}
    \item Não altera diretamente o cálculo numérico do fluxo de potência.
    \item Pode alterar atributos gráficos do diagrama.
    \item A interpretação dos resultados deve usar as tabelas e relatórios do fluxo, não a aparência do desenho.
  \end{itemize}
\end{frame}

\begin{frame}{Uso do Diagrama na Apresentação}
  O diagrama unifilar deve ser usado como apoio visual para localizar barras, ramos, geração, cargas e o ponto de falta.

  \begin{block}{Como apresentar}
    Primeiro mostrar o diagrama para situar o sistema. Depois apresentar os resultados numéricos de tensão, ângulo, fluxo, perdas e curto-circuito.
  \end{block}

  \begin{itemize}
    \item No estudo ZIP, o diagrama ajuda a identificar os ramos mais afetados, como 6--31, 8--9, 9--39 e 1--2.
    \item No estudo de falta, o diagrama ajuda a localizar a barra 4 e explicar por que barras próximas sofrem maior afundamento de tensão.
    \item O aviso de grupo base não invalida os resultados numéricos; ele afeta a importação gráfica do diagrama.
  \end{itemize}
\end{frame}

\begin{frame}{Fala Sugerida Sobre o Diagrama}
  \small
  ``O diagrama foi utilizado para visualizar a topologia do sistema e localizar os pontos analisados. Nos casos ZIP, apareceu um aviso de diferença no grupo base do diagrama, indicando inconsistência gráfica entre a tensão do diagrama e a tensão registrada no caso. Esse aviso não altera os resultados do fluxo de potência, pois os valores analisados foram retirados dos relatórios numéricos do ANAREDE. Assim, o diagrama serve como referência visual, enquanto a análise técnica é feita pelas tensões, ângulos, fluxos, perdas e correntes calculadas.''
\end{frame}

\section{Resultados Globais}

\begin{frame}{Resumo Global dos Casos}
  \begin{table}
    \centering
    \small
    \begin{tabular}{lrrrr}
      \toprule
      Caso & $P_L$ (MW) & $Q_L$ (MVAr) & $P_G$ (MW) & Perdas (MW) \\
      \midrule
      Base & 6097,1 & 1406,0 & 6140,7 & 43,61 \\
      ZIP & 6270,3 & 1433,9 & 6312,3 & 42,13 \\
      ZIP 110\% & 5768,9 & 1532,4 & 5822,0 & 53,04 \\
      \bottomrule
    \end{tabular}
  \end{table}

  \vspace{0.15cm}
  \begin{block}{Leitura dos números}
    O caso ZIP elevou a carga ativa efetiva em relação ao caso base. No cenário ZIP 110\%, a potência reativa efetiva e as perdas aumentaram, mas a carga ativa total extraída ficou menor nos resultados disponíveis.
  \end{block}
\end{frame}

\begin{frame}{Faixa de Tensões e Ângulos}
  \begin{table}
    \centering
    \small
    \begin{tabular}{lrrrr}
      \toprule
      Caso & $V_{\min}$ (\pu) & $V_{\max}$ (\pu) & $\theta_{\min}$ & $\theta_{\max}$ \\
      \midrule
      Base & 0,982 & 1,064 & -10,6$^\circ$ & 8,3$^\circ$ \\
      ZIP & 0,982 & 1,064 & -14,5$^\circ$ & 3,6$^\circ$ \\
      ZIP 110\% & 0,964 & 1,064 & -4,9$^\circ$ & 13,9$^\circ$ \\
      \bottomrule
    \end{tabular}
  \end{table}

  \vspace{0.15cm}
  \begin{itemize}
    \item O caso ZIP manteve praticamente a mesma faixa de tensão do caso base.
    \item O ZIP 110\% apresentou menor tensão mínima: 0,964 \pu{} na barra 39.
    \item Os ângulos mudaram de forma mais expressiva que as tensões.
  \end{itemize}
\end{frame}

\section{Comparação Base x ZIP}

\begin{frame}{Base x ZIP: Tensão nas Barras}
  A alteração do modelo convencional para o modelo ZIP gerou variações pequenas no módulo das tensões.

  \begin{table}
    \centering
    \small
    \begin{tabular}{lrrr}
      \toprule
      Barra & $V_{\text{Base}}$ & $V_{\text{ZIP}}$ & $\Delta V$ \\
      \midrule
      5 & 1,006 & 1,003 & -0,003 \\
      6 & 1,008 & 1,006 & -0,002 \\
      7 & 0,997 & 0,995 & -0,002 \\
      8 & 0,996 & 0,994 & -0,002 \\
      11 & 1,013 & 1,011 & -0,002 \\
      \bottomrule
    \end{tabular}
  \end{table}

  \begin{block}{Interpretação}
    O perfil de tensão permaneceu estável. A maior diferença foi de apenas 0,003 \pu, indicando que a rede manteve boa regulação de tensão nesse cenário.
  \end{block}
\end{frame}

\begin{frame}{Base x ZIP: Ângulos das Barras}
  Apesar das pequenas mudanças de tensão, os ângulos apresentaram deslocamentos maiores.

  \begin{table}
    \centering
    \small
    \begin{tabular}{lrrr}
      \toprule
      Barra & $\theta_{\text{Base}}$ & $\theta_{\text{ZIP}}$ & $\Delta\theta$ \\
      \midrule
      29 & 0,7$^\circ$ & -4,8$^\circ$ & -5,5$^\circ$ \\
      28 & -2,1$^\circ$ & -7,5$^\circ$ & -5,4$^\circ$ \\
      38 & 7,7$^\circ$ & 2,3$^\circ$ & -5,4$^\circ$ \\
      26 & -5,5$^\circ$ & -10,6$^\circ$ & -5,1$^\circ$ \\
      27 & -7,5$^\circ$ & -12,4$^\circ$ & -4,9$^\circ$ \\
      \bottomrule
    \end{tabular}
  \end{table}

  \begin{block}{Interpretação}
    A alteração angular mostra que a rede precisou redistribuir os fluxos de potência ativa para atender o novo comportamento das cargas.
  \end{block}
\end{frame}

\begin{frame}{Base x ZIP: Fluxos Ativos nos Ramos}
  Os maiores impactos apareceram nos fluxos ativos dos ramos.

  \begin{table}
    \centering
    \small
    \begin{tabular}{lrrr}
      \toprule
      Ramo & $P_{\text{Base}}$ (MW) & $P_{\text{ZIP}}$ (MW) & $\Delta P$ (MW) \\
      \midrule
      6--31 & -511,5 & -683,3 & -171,8 \\
      5--6 & -454,3 & -546,5 & -92,2 \\
      4--5 & -136,9 & -216,8 & -79,9 \\
      3--4 & 93,0 & 22,4 & -70,6 \\
      14--15 & 5,1 & 65,9 & 60,8 \\
      \bottomrule
    \end{tabular}
  \end{table}

  \begin{block}{Interpretação}
    Mesmo com pequena variação de tensão, a composição ZIP modificou a carga efetiva por barra e alterou a circulação de potência ativa.
  \end{block}
\end{frame}

\begin{frame}{Base x ZIP: Perdas}
  A perda ativa total passou de:

  \[
    43{,}61\ \text{MW} \quad \rightarrow \quad 42{,}13\ \text{MW}
  \]

  Portanto:

  \[
    \Delta P_{\text{perdas}} = -1{,}48\ \text{MW}
  \]

  \begin{block}{Interpretação}
    No caso ZIP, a redistribuição dos fluxos reduziu ligeiramente as perdas totais, apesar do aumento da carga ativa total efetiva.
  \end{block}

  \vspace{0.15cm}
  Isso mostra que perdas não dependem apenas da carga total, mas também de onde a carga está localizada e por quais caminhos a potência circula.
\end{frame}

\begin{frame}{Base x ZIP: Conclusão Parcial}
  \begin{itemize}
    \item A troca para o modelo ZIP quase não alterou a faixa de tensão.
    \item Os ângulos mudaram de forma mais relevante.
    \item Os fluxos ativos foram redistribuídos em diversos ramos.
    \item As perdas totais tiveram leve redução.
  \end{itemize}

  \begin{block}{Síntese}
    O modelo ZIP não causou problema de tensão no sistema, mas alterou a forma como a potência se distribui pela rede.
  \end{block}
\end{frame}

\section{Comparação ZIP x ZIP 110\%}

\begin{frame}{ZIP x ZIP 110\%: Tensão nas Barras}
  No cenário ZIP 110\%, apareceram quedas de tensão mais significativas.

  \begin{table}
    \centering
    \small
    \begin{tabular}{lrrr}
      \toprule
      Barra & $V_{\text{ZIP}}$ & $V_{\text{ZIP 110\%}}$ & $\Delta V$ \\
      \midrule
      39 & 1,030 & 0,964 & -0,066 \\
      9 & 1,027 & 0,968 & -0,059 \\
      1 & 1,048 & 1,000 & -0,048 \\
      8 & 0,994 & 0,969 & -0,025 \\
      7 & 0,995 & 0,973 & -0,022 \\
      \bottomrule
    \end{tabular}
  \end{table}

  \begin{block}{Interpretação}
    A barra 39 se tornou o ponto mais crítico em tensão, atingindo 0,964 \pu.
  \end{block}
\end{frame}

\begin{frame}{ZIP x ZIP 110\%: Ângulos}
  A alteração do cenário ZIP também deslocou fortemente os ângulos.

  \begin{table}
    \centering
    \small
    \begin{tabular}{lrrr}
      \toprule
      Barra & $\theta_{\text{ZIP}}$ & $\theta_{\text{ZIP 110\%}}$ & $\Delta\theta$ \\
      \midrule
      39 & -14,5$^\circ$ & 13,9$^\circ$ & 28,4$^\circ$ \\
      1 & -12,9$^\circ$ & 9,2$^\circ$ & 22,1$^\circ$ \\
      9 & -14,2$^\circ$ & 6,5$^\circ$ & 20,7$^\circ$ \\
      30 & -7,9$^\circ$ & 4,3$^\circ$ & 12,2$^\circ$ \\
      2 & -10,3$^\circ$ & 1,8$^\circ$ & 12,1$^\circ$ \\
      \bottomrule
    \end{tabular}
  \end{table}

  \begin{block}{Interpretação}
    Grandes variações angulares indicam mudança expressiva no balanço de potência ativa e no sentido de circulação de potência em parte da rede.
  \end{block}
\end{frame}

\begin{frame}{ZIP x ZIP 110\%: Fluxos Ativos}
  Os maiores desvios de fluxo ativo foram:

  \begin{table}
    \centering
    \small
    \begin{tabular}{lrrr}
      \toprule
      Ramo & $P_{\text{ZIP}}$ (MW) & $P_{\text{ZIP 110\%}}$ (MW) & $\Delta P$ (MW) \\
      \midrule
      9--39 & 23,2 & -477,5 & -500,7 \\
      8--9 & 23,4 & -471,9 & -495,3 \\
      6--31 & -683,3 & -192,1 & 491,2 \\
      1--2 & -121,1 & 313,4 & 434,5 \\
      1--39 & 121,1 & -313,4 & -434,5 \\
      \bottomrule
    \end{tabular}
  \end{table}

  \begin{block}{Interpretação}
    Há inversão ou forte alteração de sentido em ramos importantes, principalmente próximos às barras 1, 9 e 39.
  \end{block}
\end{frame}

\begin{frame}{ZIP x ZIP 110\%: Fluxos Reativos}
  A potência reativa também sofreu alterações relevantes.

  \begin{table}
    \centering
    \small
    \begin{tabular}{lrrr}
      \toprule
      Ramo & $Q_{\text{ZIP}}$ (MVAr) & $Q_{\text{ZIP 110\%}}$ (MVAr) & $\Delta Q$ (MVAr) \\
      \midrule
      8--9 & -110,9 & 60,1 & 171,0 \\
      1--39 & 30,2 & 131,7 & 101,5 \\
      1--2 & -30,1 & -131,7 & -101,6 \\
      6--31 & -94,6 & -204,7 & -110,1 \\
      2--25 & 72,5 & -14,8 & -87,3 \\
      \bottomrule
    \end{tabular}
  \end{table}

  \begin{block}{Interpretação}
    A maior circulação de reativos ajuda a explicar as quedas de tensão observadas, pois $Q$ tem relação direta com suporte e perfil de tensão.
  \end{block}
\end{frame}

\begin{frame}{ZIP x ZIP 110\%: Perdas}
  A perda ativa total passou de:

  \[
    42{,}13\ \text{MW} \quad \rightarrow \quad 53{,}04\ \text{MW}
  \]

  Portanto:

  \[
    \Delta P_{\text{perdas}} = +10{,}91\ \text{MW}
  \]

  \begin{block}{Interpretação}
    O aumento das perdas indica maior esforço elétrico nos ramos. Mesmo quando a carga ativa total efetiva não cresce, a redistribuição dos fluxos pode elevar correntes e perdas.
  \end{block}
\end{frame}

\begin{frame}{ZIP x ZIP 110\%: Ramos com Maior Variação de Perdas}
  \begin{table}
    \centering
    \small
    \begin{tabular}{lrrr}
      \toprule
      Ramo & Perda ZIP (MW) & Perda ZIP 110\% (MW) & $\Delta$ perda (MW) \\
      \midrule
      8--9 & 0,21 & 5,60 & 5,39 \\
      2--3 & 1,46 & 4,88 & 3,42 \\
      1--2 & 0,47 & 3,77 & 3,30 \\
      9--39 & 0,01 & 2,48 & 2,47 \\
      25--26 & 0,24 & 1,37 & 1,13 \\
      \bottomrule
    \end{tabular}
  \end{table}

  \begin{block}{Interpretação}
    As maiores elevações de perdas aparecem nos ramos associados às maiores mudanças de fluxo, especialmente 8--9, 2--3, 1--2 e 9--39.
  \end{block}
\end{frame}

\begin{frame}{ZIP 110\%: Conclusão Parcial}
  O cenário ZIP 110\% foi o mais severo para o sistema.

  \begin{itemize}
    \item A tensão mínima caiu para 0,964 \pu{} na barra 39.
    \item Houve grande deslocamento angular nas barras 39, 1 e 9.
    \item Alguns ramos tiveram inversão ou forte mudança no fluxo ativo.
    \item As perdas aumentaram 10,91 MW.
    \item A potência reativa total de carga subiu para 1532,4 MVAr.
  \end{itemize}

  \begin{block}{Síntese}
    O caso evidencia que mudanças no modelo e no nível de carga ZIP podem alterar bastante a operação do sistema, principalmente em tensão, ângulo, reativos e perdas.
  \end{block}
\end{frame}

\section{Discussão dos Parâmetros}

\begin{frame}{O que Acontece com a Tensão?}
  \begin{itemize}
    \item No \casobase, as tensões ficaram entre 0,982 e 1,064 \pu.
    \item No \casozip, a faixa continuou entre 0,982 e 1,064 \pu.
    \item No \casozipcento, a tensão mínima caiu para 0,964 \pu.
  \end{itemize}

  \begin{block}{Explicação pela teoria}
    A tensão é afetada pelo balanço local de potência ativa e reativa. Como o modelo ZIP altera a potência consumida conforme a tensão, ele muda o ponto de operação calculado pelo fluxo.
  \end{block}

  O efeito mais crítico apareceu quando o cenário ZIP alterado elevou o esforço reativo do sistema.
\end{frame}

\begin{frame}{O que Acontece com os Ângulos?}
  Os ângulos são fortemente relacionados ao fluxo de potência ativa.

  \[
    P_{km} \approx \frac{V_k V_m}{X_{km}}\sin(\theta_k-\theta_m)
  \]

  \begin{itemize}
    \item Quando a distribuição de carga muda, o balanço de potência ativa também muda.
    \item Para atender esse novo balanço, os ângulos das barras se ajustam.
    \item Por isso os ângulos variaram mais que os módulos de tensão.
  \end{itemize}

  \begin{block}{Resultado observado}
    A maior variação angular foi de 28,4$^\circ$ na barra 39 entre ZIP e ZIP 110\%.
  \end{block}
\end{frame}

\begin{frame}{O que Acontece com os Fluxos Ativos?}
  Os fluxos ativos mudam porque o sistema redistribui a potência entre geração e carga.

  \begin{itemize}
    \item No Base x ZIP, o maior desvio foi no ramo 6--31: -171,8 MW.
    \item No ZIP x ZIP 110\%, o maior desvio foi no ramo 9--39: -500,7 MW.
    \item Ramos próximos às barras com maiores alterações angulares foram os mais impactados.
  \end{itemize}

  \begin{block}{Interpretação operacional}
    A mudança dos fluxos pode alterar carregamentos, margens operativas e necessidade de controle de geração ou compensação.
  \end{block}
\end{frame}

\begin{frame}{O que Acontece com os Fluxos Reativos?}
  A potência reativa está diretamente associada ao suporte de tensão.

  \begin{itemize}
    \item No ZIP 110\%, a carga reativa total aumentou para 1532,4 MVAr.
    \item Houve grandes alterações de $Q$ em ramos como 8--9, 1--39, 1--2 e 6--31.
    \item As barras 39, 9 e 1 apresentaram as maiores quedas de tensão.
  \end{itemize}

  \begin{block}{Explicação pela teoria}
    Quando a demanda reativa cresce ou se redistribui, as quedas de tensão nos ramos se intensificam, principalmente em áreas com maior circulação de potência.
  \end{block}
\end{frame}

\begin{frame}{O que Acontece com as Perdas?}
  As perdas ativas nos ramos dependem da corrente:

  \[
    P_{\text{perdas}} \approx R I^2
  \]

  \begin{table}
    \centering
    \small
    \begin{tabular}{lrr}
      \toprule
      Comparação & Perda inicial & Variação \\
      \midrule
      Base $\rightarrow$ ZIP & 43,61 MW & -1,48 MW \\
      ZIP $\rightarrow$ ZIP 110\% & 42,13 MW & +10,91 MW \\
      \bottomrule
    \end{tabular}
  \end{table}

  \begin{block}{Interpretação}
    As perdas não dependem só da carga total. Elas dependem da corrente em cada ramo, da impedância e dos caminhos assumidos pelo fluxo de potência.
  \end{block}
\end{frame}

\begin{frame}{Relação entre Teoria e Resultado}
  A teoria prevê que a carga ZIP muda com a tensão:

  \[
    P,Q = f(V)
  \]

  Os resultados confirmaram essa ideia:

  \begin{itemize}
    \item as cargas efetivas mudaram entre os cenários;
    \item os ângulos se ajustaram ao novo balanço de potência ativa;
    \item os fluxos nos ramos foram redistribuídos;
    \item as perdas aumentaram ou diminuíram conforme a nova circulação de corrente;
    \item a potência reativa influenciou diretamente o perfil de tensão.
  \end{itemize}
\end{frame}

\section{Tipos de Falta}

\begin{frame}{Por que Analisar Faltas?}
  Depois do fluxo de potência, o estudo de curto-circuito avalia o comportamento do sistema diante de defeitos.

  \begin{block}{Objetivo do estudo de faltas}
    Determinar as correntes de curto-circuito, as tensões remanescentes nas barras e os impactos sobre a rede.
  \end{block}

  Durante uma falta:

  \begin{itemize}
    \item a impedância vista no ponto de defeito diminui;
    \item a corrente cresce rapidamente;
    \item a tensão no ponto de falta sofre queda acentuada;
    \item barras próximas também sofrem afundamento de tensão;
    \item os ângulos das tensões podem se deslocar, principalmente em faltas desequilibradas.
  \end{itemize}
\end{frame}

\begin{frame}{Tipos de Falta Considerados}
  Os principais tipos de falta analisados em sistemas trifásicos são:

  \begin{table}
    \centering
    \small
    \begin{tabular}{lll}
      \toprule
      Sigla & Tipo de falta & Característica \\
      \midrule
      FT & Fase-terra & Uma fase em contato com a terra \\
      FF & Fase-fase & Duas fases em curto entre si \\
      FFT & Fase-fase-terra & Duas fases em curto com a terra \\
      3F & Trifásica & Três fases em curto \\
      \bottomrule
    \end{tabular}
  \end{table}

  \begin{block}{Diferença principal}
    A falta trifásica é equilibrada. As faltas FT, FF e FFT são desequilibradas e exigem análise por componentes simétricas.
  \end{block}
\end{frame}

\begin{frame}{Componentes Simétricas}
  Para estudar faltas desequilibradas, usa-se a decomposição em três redes de sequência:

  \begin{itemize}
    \item \textbf{sequência positiva}: representa o sistema trifásico equilibrado no sentido normal de rotação;
    \item \textbf{sequência negativa}: representa desequilíbrios com sequência de fases oposta;
    \item \textbf{sequência zero}: representa correntes iguais nas três fases, normalmente associadas ao caminho pela terra ou neutro.
  \end{itemize}

  \begin{block}{Ideia central}
    Cada tipo de falta conecta as redes de sequência de uma forma diferente. Por isso as correntes, tensões e ângulos mudam conforme o curto.
  \end{block}
\end{frame}

\begin{frame}{Falta Trifásica}
  A falta trifásica é o curto-circuito simultâneo das três fases.

  \begin{itemize}
    \item É uma falta \textbf{simétrica}.
    \item As três fases permanecem equilibradas entre si.
    \item A corrente de curto costuma ser elevada.
    \item A tensão no ponto da falta tende a cair fortemente nas três fases.
    \item Os ângulos entre fases permanecem próximos de 120$^\circ$.
  \end{itemize}

  \begin{block}{Interpretação}
    Como o sistema continua equilibrado, usa-se principalmente a rede de sequência positiva. É o caso mais direto para avaliar o nível máximo de curto em muitas barras.
  \end{block}
\end{frame}

\begin{frame}{Circuito Equivalente: Falta Trifásica}
  Na falta trifásica equilibrada, apenas a rede de sequência positiva é utilizada.

  \[
    I_{3F} = \frac{V_{\text{prefalta}}}{Z_1 + Z_f}
  \]

  \begin{center}
    \begin{tabular}{c}
      \fbox{\parbox{0.72\textwidth}{
      \centering
      Fonte de sequência positiva $V_{\text{prefalta}}$
      \quad -- \quad
      impedância $Z_1$
      \quad -- \quad
      impedância de falta $Z_f$
      }}
    \end{tabular}
  \end{center}

  \begin{block}{Característica essencial}
    É uma falta simétrica: as três fases têm correntes iguais em módulo e defasadas de 120$^\circ$. Não aparecem sequências negativa e zero.
  \end{block}
\end{frame}

\begin{frame}{Falta Fase-Terra}
  A falta fase-terra ocorre quando uma fase entra em contato com a terra.

  \begin{itemize}
    \item É uma falta \textbf{desequilibrada}.
    \item A fase em falta tem grande queda de tensão.
    \item As fases sãs podem permanecer com tensão elevada ou sofrer deslocamento angular.
    \item A corrente depende fortemente do aterramento e da impedância de sequência zero.
    \item Aparecem componentes de sequência positiva, negativa e zero.
  \end{itemize}

  \begin{block}{Efeito típico}
    A tensão da fase defeituosa vai para valor muito baixo no ponto da falta, enquanto as outras fases podem ficar acima de 1,0 \pu{} dependendo da rede.
  \end{block}
\end{frame}

\begin{frame}{Circuito Equivalente: Falta Fase-Terra}
  Na falta fase-terra, as três redes de sequência ficam ligadas em série.

  \[
    I_{FT} = \frac{3V_{\text{prefalta}}}{Z_1 + Z_2 + Z_0 + 3Z_f}
  \]

  \begin{center}
    \fbox{\parbox{0.78\textwidth}{
    \centering
    $V_{\text{prefalta}}$
    \quad -- \quad
    $Z_1$
    \quad -- \quad
    $Z_2$
    \quad -- \quad
    $Z_0$
    \quad -- \quad
    $3Z_f$
    }}
  \end{center}

  \begin{block}{Característica essencial}
    Precisa de caminho de sequência zero. Por isso depende muito do aterramento de transformadores, geradores e neutros.
  \end{block}
\end{frame}

\begin{frame}{Falta Fase-Fase}
  A falta fase-fase ocorre quando duas fases entram em curto entre si, sem contato direto com a terra.

  \begin{itemize}
    \item É uma falta \textbf{desequilibrada}.
    \item As duas fases envolvidas sofrem forte perturbação.
    \item A fase não envolvida é menos afetada.
    \item Não há componente de sequência zero se não existir caminho pela terra.
    \item A corrente circula principalmente entre as duas fases em falta.
  \end{itemize}

  \begin{block}{Efeito típico}
    As tensões das fases em falta se aproximam uma da outra, reduzindo a tensão entre elas; os ângulos deixam de apresentar equilíbrio ideal.
  \end{block}
\end{frame}

\begin{frame}{Circuito Equivalente: Falta Fase-Fase}
  Na falta fase-fase, são usadas as redes de sequência positiva e negativa.

  \[
    I_{FF} = \frac{\sqrt{3}V_{\text{prefalta}}}{Z_1 + Z_2 + Z_f}
  \]

  \begin{center}
    \fbox{\parbox{0.78\textwidth}{
    \centering
    Rede de sequência positiva $Z_1$
    \quad conectada à \quad
    rede de sequência negativa $Z_2$
    \quad no ponto de falta
    }}
  \end{center}

  \begin{block}{Característica essencial}
    Não há corrente de terra idealmente. Assim, a rede de sequência zero não participa quando a falta é apenas entre duas fases.
  \end{block}
\end{frame}

\begin{frame}{Falta Fase-Fase-Terra}
  A falta fase-fase-terra ocorre quando duas fases entram em curto entre si e também com a terra.

  \begin{itemize}
    \item É uma falta \textbf{desequilibrada}.
    \item Envolve corrente entre fases e corrente para a terra.
    \item Usa redes de sequência positiva, negativa e zero.
    \item Normalmente é mais severa que a falta fase-fase.
    \item Pode produzir correntes elevadas e forte afundamento de tensão.
  \end{itemize}

  \begin{block}{Efeito típico}
    As duas fases em falta sofrem queda intensa de tensão; a fase sã também pode ser afetada por deslocamento de neutro e acoplamento da rede.
  \end{block}
\end{frame}

\begin{frame}{Circuito Equivalente: Falta Fase-Fase-Terra}
  Na falta fase-fase-terra, as três redes de sequência participam, mas a conexão equivalente é diferente da falta fase-terra.

  \begin{center}
    \fbox{\parbox{0.82\textwidth}{
    \centering
    Rede de sequência positiva $Z_1$ alimenta o ponto de falta. \\
    As redes de sequência negativa $Z_2$ e zero $Z_0$ ficam associadas ao retorno pelo defeito e pela terra.
    }}
  \end{center}

  Para falta sólida, a corrente depende de $Z_1$, $Z_2$ e $Z_0$ combinadas no ponto de falta.

  \begin{block}{Característica essencial}
    É uma falta assimétrica com caminho para terra. Costuma ser mais severa que FF porque inclui corrente de sequência zero.
  \end{block}
\end{frame}

\begin{frame}{Resumo dos Efeitos nas Grandezas}
  \begin{table}
    \centering
    \scriptsize
    \begin{tabular}{llll}
      \toprule
      Falta & Tensões & Correntes & Ângulos \\
      \midrule
      FT & Uma fase cai muito & Corrente para terra & Forte desequilíbrio \\
      FF & Duas fases perturbadas & Corrente entre fases & Perda de simetria \\
      FFT & Duas fases caem muito & Corrente entre fases e terra & Forte desequilíbrio \\
      3F & Três fases caem juntas & Corrente trifásica elevada & Mantém equilíbrio relativo \\
      \bottomrule
    \end{tabular}
  \end{table}

  \begin{block}{Ideia para apresentar}
    A falta trifásica afeta as três fases de forma equilibrada. As faltas assimétricas afetam as fases de maneira diferente, por isso aparecem tensões e ângulos desbalanceados.
  \end{block}
\end{frame}

\begin{frame}{Resumo dos Circuitos Equivalentes}
  \begin{table}
    \centering
    \scriptsize
    \begin{tabular}{llll}
      \toprule
      Falta & Redes usadas & Ligação equivalente & Característica essencial \\
      \midrule
      3F & $Z_1$ & Só sequência positiva & Curto equilibrado \\
      FT & $Z_1$, $Z_2$, $Z_0$ & Redes em série & Depende do aterramento \\
      FF & $Z_1$, $Z_2$ & Positiva e negativa & Sem sequência zero ideal \\
      FFT & $Z_1$, $Z_2$, $Z_0$ & Redes combinadas no defeito & Desequilíbrio com terra \\
      \bottomrule
    \end{tabular}
  \end{table}

  \begin{block}{Como explicar oralmente}
    O circuito equivalente muda porque cada tipo de curto impõe uma condição diferente entre fases e terra. Essa condição define quais sequências aparecem na análise.
  \end{block}
\end{frame}

\section{Resultados de Faltas}

\begin{frame}{Níveis de Curto-Circuito por Tipo de Falta}
  Resultados do ANAFAS para algumas barras representativas:

  \begin{table}
    \centering
    \small
    \begin{tabular}{lrrrr}
      \toprule
      Barra & FT (kA) & FF (kA) & FFT (kA) & 3F (kA) \\
      \midrule
      4 & 6,370 & 8,481 & 8,816 & 9,793 \\
      16 & 8,277 & 10,400 & 10,905 & 12,009 \\
      30 & 9,134 & 8,605 & 9,639 & 9,937 \\
      39 & 30,608 & 27,271 & 31,107 & 31,490 \\
      \bottomrule
    \end{tabular}
  \end{table}

  \begin{block}{Interpretação}
    A barra 39 apresentou os maiores níveis de curto entre os dados extraídos, chegando a 31,490 kA na falta trifásica.
  \end{block}
\end{frame}

\begin{frame}{Comparação dos Tipos de Falta}
  Observa-se que a corrente de falta depende do tipo de defeito e da barra analisada.

  \begin{itemize}
    \item Na barra 4, a falta trifásica foi a maior: 9,793 kA.
    \item Na barra 16, a falta trifásica também foi a maior: 12,009 kA.
    \item Na barra 39, os valores foram muito altos para todos os tipos.
    \item A falta fase-terra na barra 39 chegou a 30,608 kA, mostrando forte contribuição da rede e do aterramento.
  \end{itemize}

  \begin{block}{Leitura técnica}
    A falta trifásica tende a ser a referência para nível de curto simétrico, mas faltas envolvendo terra podem ser muito severas dependendo das impedâncias de sequência zero.
  \end{block}
\end{frame}

\begin{frame}{Falta FT na Barra 4: Tensões Durante a Falta}
  Para a falta fase-terra na barra 4, o ANAFAS indicou corrente de falta de:

  \[
    I_f = 6{,}37\ \text{kA}
  \]

  Tensões em barras próximas:

  \begin{table}
    \centering
    \scriptsize
    \begin{tabular}{lrrrrrr}
      \toprule
      Barra & $V_a$ & $\angle a$ & $V_b$ & $\angle b$ & $V_c$ & $\angle c$ \\
      \midrule
      2 & 0,766 & -0,3$^\circ$ & 1,004 & -120,6$^\circ$ & 1,007 & 120,5$^\circ$ \\
      3 & 0,511 & 0,0$^\circ$ & 1,065 & -125,9$^\circ$ & 1,070 & 125,7$^\circ$ \\
      4 & 0,000 & 0,0$^\circ$ & 1,211 & -134,5$^\circ$ & 1,215 & 134,4$^\circ$ \\
      5 & 0,360 & -1,4$^\circ$ & 1,088 & -127,0$^\circ$ & 1,084 & 127,2$^\circ$ \\
      \bottomrule
    \end{tabular}
  \end{table}
\end{frame}

\begin{frame}{Interpretação da Falta FT na Barra 4}
  A falta fase-terra foi aplicada na fase A da barra 4.

  \begin{itemize}
    \item Na barra 4, $V_a = 0{,}000$ \pu, pois a fase A está em curto para a terra.
    \item As fases B e C na barra 4 subiram para aproximadamente 1,21 \pu.
    \item Nas barras próximas, a fase A também sofreu afundamento: 0,511 \pu{} na barra 3 e 0,360 \pu{} na barra 5.
    \item Os ângulos das fases B e C se deslocaram para cerca de $\pm 134^\circ$, mostrando desequilíbrio.
  \end{itemize}

  \begin{block}{Conclusão do caso}
    O resultado confirma a teoria: em falta fase-terra, a fase defeituosa sofre queda severa de tensão, enquanto as fases sãs ficam desequilibradas em módulo e ângulo.
  \end{block}
\end{frame}

\begin{frame}{Relação entre Fluxo de Potência e Falta}
  O fluxo de potência e o curto-circuito são estudos complementares.

  \begin{itemize}
    \item O fluxo de potência determina o estado pré-falta: tensões, ângulos, gerações e carregamentos.
    \item O estudo de falta usa esse estado para avaliar correntes e tensões durante o defeito.
    \item Tensões pré-falta influenciam a corrente calculada.
    \item A localização da falta define quais barras terão maior afundamento de tensão.
  \end{itemize}

  \begin{block}{Ligação com a teoria}
    O material da disciplina destaca que, em termos práticos, antes da análise de curto-circuito é necessário executar o fluxo de potência para conhecer as tensões pré-falta.
  \end{block}
\end{frame}

\section{Conclusões}

\begin{frame}{Conclusões Principais}
  \begin{enumerate}
    \item O \casozip{} manteve tensões próximas ao \casobase, mas alterou ângulos e fluxos.
    \item A modelagem ZIP modificou a carga efetiva por barra, conforme previsto pela teoria.
    \item A comparação Base x ZIP apresentou pequena redução de perdas: -1,48 MW.
    \item O \casozipcento{} foi o cenário mais severo, com queda de tensão até 0,964 \pu{} e aumento de perdas de 10,91 MW.
    \item Nas faltas, as tensões e correntes mudaram conforme o tipo de curto-circuito e a localização do defeito.
    \item A falta FT na barra 4 confirmou o comportamento esperado: tensão da fase defeituosa igual a zero e desequilíbrio nas fases sãs.
  \end{enumerate}
\end{frame}

\begin{frame}{Conclusão Final}
  \begin{block}{Resultado final}
    A análise mostra que a modelagem ZIP é importante porque representa de forma mais realista a dependência da carga em relação à tensão.
  \end{block}

  \vspace{0.25cm}
  No sistema estudado, a mudança para ZIP não provocou grandes violações de tensão, mas alterou a distribuição dos fluxos e dos ângulos.

  \vspace{0.25cm}
  Quando o cenário ZIP foi modificado, os impactos se tornaram mais fortes: queda de tensão, aumento de potência reativa circulante, mudança de sentido em alguns fluxos e aumento das perdas.

  \vspace{0.25cm}
  Na análise de faltas, os resultados também foram coerentes com a teoria: faltas simétricas afetam as três fases de forma equilibrada, enquanto faltas assimétricas provocam desequilíbrio entre módulos e ângulos das tensões.

  \vspace{0.25cm}
  Portanto, os resultados obtidos no ANAREDE e no ANAFAS estão coerentes com a teoria apresentada no material da disciplina.
\end{frame}

\begin{frame}{Sugestão de Fala para Encerramento}
  \small
  ``Com base na teoria, o fluxo de potência determina o estado operativo do sistema por meio das tensões, ângulos, fluxos e perdas. Ao aplicar o modelo ZIP, a carga deixa de ser fixa e passa a depender da tensão em cada barra. Nos nossos resultados, isso aparece claramente: no caso ZIP, as tensões mudaram pouco, mas os ângulos e os fluxos foram redistribuídos. No cenário ZIP 110\%, a condição ficou mais severa, com maior queda de tensão, maior circulação de reativos e aumento das perdas. Na parte de faltas, a falta trifásica representa o defeito equilibrado, enquanto FT, FF e FFT provocam desequilíbrios nas tensões e nos ângulos. Assim, os resultados confirmam que tanto o modelo de carga quanto o tipo de falta influenciam diretamente a resposta do sistema elétrico.''
\end{frame}

\begin{frame}
  \centering
  \Huge Obrigado

  \vspace{0.7cm}
  \Large Perguntas?
\end{frame}

\end{document}
